\documentclass[12pt]{article}
\usepackage[utf8]{inputenc}
\usepackage[T1]{fontenc}
\usepackage{hyperref}
\title{Harsusi Language Resources for K-12 Curriculum Development}
\author{Compiled by Claude Sonnet 5.5 (claude-sonnet-5-5) for Andr\'as Kornai}
\date{}
\begin{document}
\maketitle

\section{Introduction}
Harsusi (hss) is an endangered Modern South Arabian language of the Jiddat al-Harasis, Oman, with an estimated 600 to 1,000 speakers. It is unwritten. Documentation is limited to linguistic work; no educational materials were found.

\section{Primary Educational Resources}
None found in search results.

\section{Linguistic and Reference Materials}
\begin{itemize}
\item Harry Stroomer, \emph{Harsusi Texts from Oman} (2004), publishing texts collected by T.M. Johnstone; the book notes how little documentation exists. \url{https://reado.app/en/book/harsusi-texts-from-oman-harry-stroomer/9783447050975}
\item Leeds PhD project of Hammal Al Balushi on Harsusi phonology and morphology. \url{https://ahc.leeds.ac.uk/languages/pgr/2639/hammal-al-balushi}
\end{itemize}

\section{Religious and Community Resources}
None found.

\section{Digital and Audio Resources}
\begin{itemize}
\item Endangered Languages Archive collection ``The documentation and ethnolinguistic analysis of Modern South Arabian: Harsusi'' (audio and audio-visual recordings by Domenyk Eades and Miranda Morris). \url{https://eldp.access.preservica.com/dk0314/}
\item Endangered Languages Project sample page for Harsusi. \url{https://www.endangeredlanguages.com/lang/4281/samples/9211}
\end{itemize}

\section{Curriculum Development Context}
Harsusi is not known to be taught in schools and has no writing system. Any curriculum would need community-led orthography development, drawing on the archive recordings.

\end{document}